\section{如何成为 Agent Manager}

\subsection{软件开发团队的演进}

软件开发团队的组织形态正在经历第四次重大转型：

\begin{enumerate}
    \item \textbf{1940--1960s}：单个开发者管理整个项目
    \item \textbf{1970--1990s}：软件团队主流化，专业分工出现
    \item \textbf{2023--2025}：开发者借助 AI 编程系统辅助工作
    \item \textbf{2025+}：单个开发者管理多个 AI Agent 的工作输出
\end{enumerate}

\begin{importantbox}{从管理人到管理 Agent}
开发的演化路径是：管理自己的输出 $\rightarrow$ 管理团队的输出 $\rightarrow$ 管理 AI 辅助团队的输出 $\rightarrow$ \textbf{单人管理多个 AI Agent 的输出}。最后一步是当前正在发生的变革。
\end{importantbox}

\subsection{软件任务的分解}

一个典型的软件开发任务可以分解为以下步骤，其中绿色表示人类主导，蓝色表示 Agent 可以承担：

\begin{enumerate}
    \item 提供高层需求 \textbf{{[}人类{]}}
    \item 将需求转化为设计文档 \textbf{{[}人类/Agent{]}}
    \item 根据设计文档实现方案 \textbf{{[}Agent{]}}
    \item 添加测试 \textbf{{[}Agent{]}}
    \item 确保 CI 通过 \textbf{{[}Agent{]}}
    \item 代码审查 \textbf{{[}Agent{]}}
    \item 更新文档 \textbf{{[}Agent{]}}
\end{enumerate}

\begin{knowledgebox}{人类的核心价值在第 1--2 步}
在 Agent 管理模式下，人类工程师的最高价值体现在\textbf{需求理解}和\textbf{架构设计}上。这两步需要业务上下文、技术判断力和``品味''——恰好是当前 AI 最薄弱的环节。
\end{knowledgebox}

\subsection{指导 Agent 的四种技术}

\subsubsection{Agent 行为文件}

如 CLAUDE.md、.cursorrules、AGENTS.md 等，提供项目级别的行为指导。

\subsubsection{Hooks}

确定性脚本，在预定义的事件触发时运行：
\begin{itemize}
    \item \texttt{PreToolUse}：工具调用前
    \item \texttt{PostToolUse}：工具调用后
    \item \texttt{UserPromptSubmit}：用户提交 prompt 时
    \item \texttt{PreCompact}：context 压缩前
\end{itemize}

\subsubsection{Commands}

将常用 prompt 保存为文件，Agent 可以按需执行。典型用例：运行测试、审查代码、提交 git commit 并推送。

\subsubsection{Subagents}

运行时委派机制，核心目的：
\begin{itemize}
    \item 为不同类型的工作创建独立的开发者``人格''（frontend、backend 等）
    \item 清晰分离不同工作流的 context
    \item 提供定制化的 system prompt、工具集和独立的 context window
\end{itemize}

\begin{importantbox}{Subagent 是 Agent 管理 Agent 的雏形}
Subagent 模式代表了一个重要趋势：Agent 管理其他 Agent。主 Agent 负责任务分解和委派，subagent 专注于特定领域的执行，每个 subagent 拥有独立的 context 避免信息过载。
\end{importantbox}

\subsection{最佳实践}

\begin{enumerate}
    \item \textbf{设置严格的防护栏}：代码库中的测试 + CI/CD 最佳实践是不可或缺的 backstop
    \item \textbf{审计每个 Agent 的修改}：标记（label）每个由 Agent 产生的 diff
    \item \textbf{不同任务使用不同模型}：简单任务用快速模型，复杂任务用强模型
    \item \textbf{复杂任务需要更多前期引导}：不是所有任务都适合完全异步委派
    \item \textbf{频繁 checkpoint（commit）}：定期提交代码，防止 Agent 偏离方向时丢失进展
\end{enumerate}

\begin{warningbox}{开放问题}
\begin{itemize}
    \item 如何自动化每个任务最初 10--20\% 的调研阶段？
    \item 如何维护一个待执行任务的队列？（对于一次性修改较容易，但对持续性工作流仍是挑战）
\end{itemize}
\end{warningbox}

\subsection{本章小结}

Agent Manager 模式要求开发者从``写代码的人''转型为``管理写代码的 Agent 的人''。核心技术包括行为文件、hooks、commands 和 subagents。关键原则是：设置防护栏、审计修改、匹配任务与模型、频繁 checkpoint。

